\section{Условие}
Запустить тестовую программу:
\begin{listing}[Тестовая программа]{lst:test-prog}{Prolog}
parent(don, randy).
parent(don, mike).
parent(don, anne).
parent(rosie, randy).
parent(rosie, mike).
parent(rosie, anne).

?-parent(don, mike).
\end{listing}

Проанализировать реакцию системы и множество ответов при разных аргументах в вопросе.

Разработать свою программу --- <<Телефонный справочник и автомобили>>. Абоненты могут
иметь несколько телефонов. Протестировать работу программы, используя разные вопросы.
\begin{itemize}
    \item <<Телефонный справочник>>: Фамилия, Номер телефона, Адрес --- структура (Город, Улица, Номер дома, Номер квартиры).
    \item <<Автомобили>>: Фамилия\_владельца, Марка, Цвет, Стоимость, Номер.
\end{itemize}

Владелец может иметь несколько телефонов, автомобилей (Факты). В разных городах есть 
однофамильцы, в одном городе --- фамилия уникальна. Используя конъюнктивное правило и
простой вопрос, обеспечить возможность поиска: по Марке и Цвету автомобиля найти Фамилию, Город, Телефон.